
Ziel dieser Arbeit war es, zu untersuchen, welche Voraussetzungen kleine und mittlere Unternehmen erfüllen müssen, damit 
ERP-Systeme als technologische Basis digitaler Transformation wirksam werden können. Die Analyse verdeutlicht, dass ERP-Systeme 
für KMU weit mehr als operative Standardsoftware darstellen können. Durch die Integration betrieblicher Prozesse und die 
Bereitstellung einer gemeinsamen Datenbasis schaffen sie grundsätzlich Voraussetzungen für Transparenz, bereichsübergreifende 
Zusammenarbeit und datenbasierte Entscheidungen. ERP-Systeme sind jedoch nicht mit digitaler Transformation gleichzusetzen. 
Ihre Einführung allein führt nicht automatisch zu einer nachhaltigen Veränderung von Prozessen, Entscheidungsstrukturen oder 
Geschäftsmodellen.

Die Forschungsfrage lässt sich daher dahingehend beantworten, dass ERP-Systeme insbesondere dann als technologische Basis 
digitaler Transformation wirksam werden können, wenn zwischen den Anforderungen des Systems und den organisatorischen Voraussetzungen 
des KMU eine ausreichende Passung besteht. Diese Passung wurde im Rahmen der Arbeit durch fünf Dimensionen konkretisiert: 
Prozess-, Daten-, Ressourcen-, Kompetenz- und Strategie-Fit.

Ein hoher Prozess-Fit setzt voraus, dass zentrale Geschäftsprozesse ausreichend verstanden, strukturiert und für eine Abbildung 
im ERP-System geeignet sind. Der Daten-Fit betrifft die Qualität, Konsistenz und Überführbarkeit relevanter Unternehmensdaten in 
eine gemeinsame Datenbasis. Der Ressourcen-Fit umfasst die realistische Verfügbarkeit von Budget, Zeit und personellen Kapazitäten. 
Ergänzend beschreibt der Kompetenz-Fit, ob Mitarbeitende und Führungskräfte über ausreichendes Wissen, Akzeptanz und 
Veränderungsbereitschaft verfügen. Der Strategie-Fit berücksichtigt schließlich, ob die ERP-Einführung in übergeordnete 
Unternehmens- und Digitalisierungsziele eingebettet ist.

Das entwickelte ERP-KMU-Fit-Modell bündelt diese Voraussetzungen in einem gemeinsamen Bezugsrahmen. Es unterstützt eine 
ganzheitliche Betrachtung der ERP-Einführung und verdeutlicht, dass Schwächen in einer Dimension die Wirkung anderer Dimensionen 
begrenzen können. Damit dient das Modell als Orientierungsinstrument, um Handlungsbedarfe frühzeitig sichtbar zu machen und 
ERP-Projekte nicht ausschließlich als technische, sondern als organisatorische und strategische Veränderungsvorhaben zu verstehen.

Zusammenfassend zeigt sich, dass ERP-Systeme für KMU eine wichtige Grundlage digitaler Transformation bilden können, ihr Nutzen 
jedoch vom Zusammenspiel technologischer, organisatorischer und strategischer Voraussetzungen abhängt. Nicht der Funktionsumfang 
eines Systems entscheidet über seine Transformationswirkung, sondern die Passung zwischen ERP-System, Prozessen, Daten, 
Ressourcen, Kompetenzen und strategischen Zielen des Unternehmens.
